\documentclass[11pt]{amsart}
\usepackage{amsmath,amssymb,amsthm}
\usepackage[margin=1in]{geometry}
\usepackage{enumitem}
\usepackage{url}

\theoremstyle{plain}
\newtheorem{theorem}{Theorem}[section]
\newtheorem{proposition}[theorem]{Proposition}
\newtheorem{conjecture}[theorem]{Conjecture}
\theoremstyle{definition}
\newtheorem{definition}[theorem]{Definition}
\theoremstyle{remark}
\newtheorem{remark}[theorem]{Remark}

\title[Odd greedy orbits with bounded numerators]{Odd greedy orbits with bounded numerators:\\partial results and open problems}

\author{Quanzhi Wang}
\email{mrhwrn0504@gmail.com}

\date{10 November 2026}

\begin{document}
\begin{abstract}
We study the odd greedy expansion of fractions and the orbits it generates. For each $A\ge 2$ we ask whether there exists an infinite lawful orbit all of whose reduced numerators lie in $[2,A]$. We prove that no such orbit exists for $A\le 4$ by an elementary argument, and for $A\le 15$ by a computer-assisted finite-state certificate, which we check by two independent implementations. For $A=16$ we reduce the problem to a single periodic word of period eight and show that this word is $p$-adically admissible at the primes $2$, $11$, $13$, with an explicit fourteen-step integer seed. The general statement, and the unbounded version in which numerators exceed every bound infinitely often, remain open. We state these open problems precisely and record the gaps that separate the results proved here from the full conjecture.
\end{abstract}

\maketitle

\section{Introduction and status of results}

Every result below carries a status label. The labels are used consistently throughout:
\begin{itemize}[leftmargin=2em]
\item \emph{Proved (elementary)}: proved by hand.
\item \emph{Proved (computer-assisted)}: depends on a finite computation whose method is described and whose code is to be released; the result has not yet been independently rerun by a third party.
\item \emph{Reduction, incomplete}: a reduction whose source has not yet been independently re-derived.
\item \emph{Open}: not proved and not refuted.
\end{itemize}

\begin{center}
\begin{tabular}{p{0.52\textwidth} p{0.28\textwidth} p{0.12\textwidth}}
\hline
Statement & Status & Section\\
\hline
No infinite lawful orbit with $2\le a_k\le 4$ for all $k$ & Proved (elementary; re-verification pending) & \S\ref{sec:small}\\
No infinite lawful orbit with $2\le a_k\le 15$ for all $k$ & Proved (computer-assisted) & \S\ref{sec:comp}\\
Periodic word $W$ of period $8$ is $p$-adically admissible at $2,11,13$ & Proved (computer-assisted) & \S\ref{sec:sixteen}\\
Reduction of $A=16$ to $W$ (and $A\le 24$ to $A=16$) & Reduction, incomplete & \S\ref{sec:sixteen}\\
The case $A\ge 16$ & Open & \S\ref{sec:open}\\
Unbounded numerators in the strong sense & Open & \S\ref{sec:open}\\
The full conjecture for all $A\ge 2$ & Conjecture, open & \S\ref{sec:open}\\
\hline
\end{tabular}
\end{center}

\section{Setting and definitions}\label{sec:setting}

Let $a/b$ be a reduced fraction with $b$ odd. The \emph{odd greedy step} is defined as follows. Let $q$ be the smallest odd positive integer with $aq\ge b$. Put
\[
r = aq-b,\qquad g=\gcd(r,q^2),\qquad a' = \frac{r}{g},\qquad b' = \frac{bq}{g}.
\]
Then $a'/b'$ is again a reduced fraction with odd denominator, and we iterate.

\begin{definition}
A \emph{lawful orbit} is a sequence $a_1/b_1, a_2/b_2,\dots$ generated by repeated application of the odd greedy step. It is \emph{$A$-bounded} if $2\le a_k\le A$ for all $k$.
\end{definition}

\begin{conjecture}\label{conj:main}
For every $A\ge 2$ there is no infinite $A$-bounded lawful orbit. Equivalently, every infinite lawful orbit satisfies $\sup_k a_k=\infty$.
\end{conjecture}

Finite orbits and orbits verified only numerically are not counterexamples. A counterexample must be an infinite lawful orbit.

\section{Small bounds: $A\le 4$}\label{sec:small}

\begin{proposition}[Proved, elementary; re-verification pending]
For $A\le 4$ there is no infinite $A$-bounded lawful orbit.
\end{proposition}

\begin{proof}[Proof sketch]
Write $\delta_k=r_k/a_k$. When $a_k\le A$, the quantity $\delta_k$ takes values in a finite set $D_A$, and consecutive values satisfy $|\delta_{k+1}-\delta_k|\ge 1/A^2$. The recurrence $q_{k+1}=q_k^2/\delta_k-q_k+\delta_{k+1}$ together with the finiteness of $D_A$ rules out an infinite path for $A\le 4$. The full case analysis, and the two-step identity
\[
a_k n_k^2+g_kg_{k+1}a_{k+2}=g_k a_{k+1}(n_k+n_{k+1}),
\]
are to be re-derived line by line and checked by a short script before this result is cited as proved.
\end{proof}

\section{Computer-assisted bound: $A\le 15$}\label{sec:comp}

\begin{proposition}[Proved, computer-assisted]
For $A\le 15$ there is no infinite $A$-bounded lawful orbit.
\end{proposition}

\paragraph{Method.} Encode a step as a state $(a,r,q\bmod M)$ for a suitable modulus $M$. Compute all transitions out of states with $a\le A$. Repeatedly delete states with no successor inside the set. If the surviving set is empty, no infinite bounded orbit exists in the residue-class model; a state whose lift is an infinite lawful orbit would survive every deletion round, so emptiness proves the proposition.

\paragraph{Two implementations.} The first (\texttt{trim}) works on $(a,r,q\bmod M)$. The second (\texttt{btrim}) works directly on $(a,b\bmod M)$ with integer arithmetic. With $M=720720$, \texttt{btrim} gives the empty set for $9\le A\le 15$, agreeing with \texttt{trim}.

\paragraph{Sanity check.} For $2\le A\le 8$ the exact maximal length of a bounded run, computed by residue-class depth-first search, is $0,1,2,4,5,7,7$.

\paragraph{Limitation.} The conclusion depends on the choice of $M$ and on the correctness of the implementation. It has not yet been rerun by an independent party. The code and checksums will be released with the paper.

\section{The case $A=16$ and the eight-cycle}\label{sec:sixteen}

\begin{remark}[Reduction, source to be verified]
The claim that every $A\le 24$ reduces to $A=16$, and that for $A=16$ the only candidate structure is a single eight-cycle, comes from an archived research record (rounds R49/R50). It has not been re-derived here; our independent search agrees that the eight-cycle below is the only candidate found.
\end{remark}

Let $W=(8,11,16,13,16,11,8,13)$, all of whose steps have $g=1$.

\begin{proposition}[Proved, computer-assisted]
The word $W$ is $p$-adically admissible at $p=2,11,13$, and every branch at these primes is a full branch.
\end{proposition}

\begin{proposition}[Proved, computer-assisted]
The integer $q_0=491458458962611 \bmod 876948437860352$ gives the starting fraction $8/3931667671700877$, whose orbit follows the word $W$ for fourteen steps.
\end{proposition}

\begin{remark}
These results do not establish the existence of an infinite orbit. Admissibility is checked at only three primes, and a fourteen-step seed is not an infinite orbit. The existence of a seed admissible at all primes simultaneously is not established.
\end{remark}

\section{Heuristic evidence (not a proof)}\label{sec:heur}

We sampled $6000$ fractions with $300$-bit odd denominators and numerators in $[2,200)$. Of these, $4872$ were tracked and all terminated within $400$ steps, the longest after $17$ steps. The remaining $1128$ samples are not included in the statistics; the reason for their exclusion still needs to be checked. The fraction of samples whose orbit lasts more than $n$ steps is
\begin{center}
\begin{tabular}{rr}
\hline
$n$ & $\Pr(\text{length}>n)$\\
\hline
3 & $0.693$\\
6 & $0.216$\\
9 & $0.0394$\\
12 & $0.00554$\\
15 & $0.00041$\\
16 & $0.00021$\\
\hline
\end{tabular}
\end{center}
For the family $3/(N!-1)$ with $N\in\{5,\dots,60,80,100,150,200,300,400\}$ every orbit terminates, with an ascent of length about $N$; $N=49$ is unresolved due to loss of precision. These data are consistent with the conjecture but are not evidence that could be promoted to a theorem.

\section{Open problems}\label{sec:open}

\begin{conjecture}[Open]
Conjecture~\ref{conj:main} holds for all $A\ge 2$.
\end{conjecture}

\begin{conjecture}[Open, unbounded version]\label{conj:unbounded}
Every infinite lawful orbit has $\sup_k a_k=\infty$, and moreover for every $M$ there are infinitely many $k$ with $a_k>M$.
\end{conjecture}

\begin{remark}
The three gaps are: (i) the case $A\ge 16$ requires a single integer seed admissible at all primes, (ii) the unbounded version requires a closing argument, and (iii) the density argument has a gap at a critical root. All three reduce to the same difficulty: passing from ``every finite length has a seed'' to ``one seed works for all lengths.'' No known method achieves this step.
\end{remark}

\section{Discussion and next steps}

The three gaps identified above reduce to two statements. The version of the unbounded problem in which, for every $M$, infinitely many $k$ satisfy $a_k>M$ is equivalent to Conjecture~\ref{conj:main}: an infinite orbit that is eventually bounded by $M$ restarts, from the first index after which it stays bounded, as an infinite $M$-bounded orbit. Hence the case $A\ge 16$ is subsumed by the unbounded problem and need not be treated separately. The density result is a different and weaker statement.

\subsection*{Route survey and verification status}

\begin{center}
\begin{tabular}{p{0.26\textwidth} p{0.30\textwidth} p{0.36\textwidth}}
\hline
Route & Uniform in $A$? & Verification status\\
\hline
Subgreedy finiteness (Eppstein): every fraction has finitely many subgreedy expansions & Yes & Computationally supported for $y\le 2500$. It implies termination of odd greedy only through a distinctness lemma for switch-point expansions, which is not proved. Naive switching gives repeated expansions, as expected when the ordinary greedy tail coincides with the odd continuation.\\
Nested residue-class obstruction (Mahler type, cf.\ Kumar) & No, only for fixed $A$ & Fails for the eight-cycle $W$: the existence-type survival set is nonempty at every prime $p\le 47$ tested, so no single-prime obstruction is available.\\
Density (Collatz type, cf.\ Tao) & No, gives almost all starting points & Not a route to the conjecture; a different target.\\
Explicit long constructions (Brown, Wagon) & No & Empirically the family $3/(N!-1)$ terminates for all $N$ tested up to $400$, except $N=49$ where precision was lost.\\
\hline
\end{tabular}
\end{center}

\noindent The literature found does not contain a proof of any uniform route. The counting results of Conlon et al.\ and the Graham–Butler results concern different problems and do not apply directly.

\subsection*{Status}

The full conjecture remains open. The only route found that is uniform in $A$ is subgreedy finiteness, which is itself a conjecture and requires a distinctness lemma that has not been proved. Before submission the following must be completed: a proof or refutation of the distinctness lemma; an independent rerun of the $A\le 15$ certificate; and verification of every cited reference against its original source.

\begin{thebibliography}{99}
\bibitem{erdos} Erd\H{o}s Problems, Problem \#282. \url{https://www.erdosproblems.com} (status to be checked on the day of submission).
\bibitem{pihko} P. Pihko et al., work on greedy expansions of fractions. [to be checked]
\bibitem{lm} Louwsma and Martino, on odd greedy expansions. [to be checked]
\bibitem{koizumi} Koizumi, on related greedy expansions. [to be checked]
\bibitem{eppstein} D. Eppstein, subgreedy conjecture. [to be checked]
\bibitem{brown} K. S. Brown, long-run constructions $3/179$, $5/5809$, $3/(N!-1)$. [to be checked]
\bibitem{tao} T. Tao, almost all Collatz orbits (density method, for comparison). [to be checked]
\end{thebibliography}

\end{document}
